\begin{frame}
    \frametitle{Ejercicio}

    \begin{exercise}
        Calcula la dimensión de los siguientes espacios vectoriales
        reales y da una base de cada uno:
        \begin{enumerate}
            \item $\mathbb{R}^3\times\mathbb{R}^2$;
            \item $M_{2\times3}(\mathbb{R})\times\mathbb{R}_3[x]$;
            \item $\mathbb{R}_1[x]\times\mathbb{R}_1[x]\times\mathbb{R}_1[x]$.
        \end{enumerate}
    \end{exercise}
\end{frame}


\begin{frame}[allowframebreaks]
    \frametitle{Ejercicio}

    \begin{solution}
        Aplicamos el teorema sobre bases del producto y su corolario.

        \begin{enumerate}
            \item $\dim=3+2=5$. Una base es
                  \[
                      \{(e_1,0),(e_2,0),(e_3,0),(0,e'_1),(0,e'_2)\},
                  \]
                  donde $e_i$ y $e'_j$ son los vectores canónicos de
                  $\mathbb{R}^3$ y $\mathbb{R}^2$. Con la identificación
                  $\mathbb{R}^3\times\mathbb{R}^2=\mathbb{R}^5$, es la base
                  canónica de $\mathbb{R}^5$.

            \item $\dim=6+4=10$. Una base está formada por los seis pares
                  $(E_{ij},0)$, con $i=1,2$ y $j=1,2,3$, y los cuatro pares
                  $(0,1),(0,x),(0,x^2),(0,x^3)$.

            \item $\dim=2+2+2=6$. Una base es
                  \[
                      (1,0,0),\ (x,0,0),\ (0,1,0),\ (0,x,0),\ (0,0,1),\ (0,0,x).\qedhere
                  \]
        \end{enumerate}
    \end{solution}
\end{frame}


\begin{frame}
    \frametitle{Ejercicio}

    \begin{exercise}
        En $\mathbb{R}^2\times\mathbb{R}_1[x]$ se consideran los vectores
        \[
            u_1=\bigl((1,0),\,1+x\bigr),
            \qquad
            u_2=\bigl((0,1),\,x\bigr),
        \]
        \[
            u_3=\bigl((1,1),\,0\bigr),
            \qquad
            u_4=\bigl((0,0),\,1\bigr).
        \]
        Decide si $\{u_1,u_2,u_3,u_4\}$ es una base de
        $\mathbb{R}^2\times\mathbb{R}_1[x]$.
    \end{exercise}

    \medskip

    \alert{Indicación:} trabaja con coordenadas respecto de la base
    \[
        B=\bigl\{((1,0),0),\ ((0,1),0),\ ((0,0),1),\ ((0,0),x)\bigr\}.
    \]
\end{frame}


\begin{frame}[allowframebreaks]
    \frametitle{Ejercicio}

    \begin{solution}
        Los vectores de coordenadas en $B$ son
        \[
            (u_1)_B=\begin{pmatrix}1\\0\\1\\1\end{pmatrix},\quad
            (u_2)_B=\begin{pmatrix}0\\1\\0\\1\end{pmatrix},\quad
            (u_3)_B=\begin{pmatrix}1\\1\\0\\0\end{pmatrix},\quad
            (u_4)_B=\begin{pmatrix}0\\0\\1\\0\end{pmatrix}.
        \]
        Sea $A$ la matriz $4\times4$ que tiene estos vectores por columnas.

        Dado $v\in\mathbb{R}^2\times\mathbb{R}_1[x]$, por las propiedades
        de la aplicación de coordenadas y la unicidad de las coordenadas,
        \[
            v=\lambda_1u_1+\cdots+\lambda_4u_4
            \iff
            v_B=\lambda_1(u_1)_B+\cdots+\lambda_4(u_4)_B
            \iff
            A\begin{pmatrix}\lambda_1\\\vdots\\\lambda_4\end{pmatrix}=v_B.
        \]

        \framebreak

        Escalonamos $A$:
        \[
            \begin{pmatrix}
                1&0&1&0\\
                0&1&1&0\\
                1&0&0&1\\
                1&1&0&0
            \end{pmatrix}
            \xrightarrow[F_4-F_1]{F_3-F_1}
            \begin{pmatrix}
                1&0&1&0\\
                0&1&1&0\\
                0&0&-1&1\\
                0&1&-1&0
            \end{pmatrix}
            \xrightarrow[F_4-2F_3]{F_4-F_2}
            \begin{pmatrix}
                1&0&1&0\\
                0&1&1&0\\
                0&0&-1&1\\
                0&0&0&-2
            \end{pmatrix}.
        \]
        (La segunda flecha indica que primero se hace $F_4-F_2$ y después
        $F_4-2F_3$.)

        Así, $\rg(A)=4$: hay un pivote en cada fila y en cada columna.
        Por tanto, para cada $v$, el sistema $A\lambda=v_B$ tiene
        \alert{solución única}, es decir, $v$ se escribe de forma única
        como combinación lineal de $u_1,u_2,u_3,u_4$.

        Por el teorema de representación única, $\{u_1,u_2,u_3,u_4\}$
        es una base de $\mathbb{R}^2\times\mathbb{R}_1[x]$.
    \end{solution}
\end{frame}


\begin{frame}
    \frametitle{Ejercicio}

    \begin{exercise}
        Sean $V$ y $W$ espacios vectoriales sobre $k$.
        \begin{enumerate}
            \item Demuestra que si $v_1,\ldots,v_r\in V$ son linealmente
                  independientes, entonces, para cualesquiera
                  $w_1,\ldots,w_r\in W$, los vectores
                  \[
                      (v_1,w_1),\ldots,(v_r,w_r)
                  \]
                  son linealmente independientes en $V\times W$.
            \item Muestra con un ejemplo que el recíproco es falso.
            \item Si $v_1,\ldots,v_r$ generan $V$ y $w_1,\ldots,w_r$
                  generan $W$, ¿generan necesariamente
                  $(v_1,w_1),\ldots,(v_r,w_r)$ el espacio $V\times W$?
        \end{enumerate}
    \end{exercise}
\end{frame}


\begin{frame}[allowframebreaks]
    \frametitle{Ejercicio}

    \begin{solution}
        \alert{(1)} Supongamos que
        $\lambda_1(v_1,w_1)+\cdots+\lambda_r(v_r,w_r)=(0,0)$.
        Calculando componente a componente,
        \[
            \bigl(\lambda_1v_1+\cdots+\lambda_rv_r,\;
                  \lambda_1w_1+\cdots+\lambda_rw_r\bigr)=(0,0).
        \]
        En particular, la primera componente da
        $\lambda_1v_1+\cdots+\lambda_rv_r=0$ en $V$, y la independencia
        lineal de $v_1,\ldots,v_r$ implica $\lambda_1=\cdots=\lambda_r=0$.
        No ha hecho falta usar la segunda componente.

        \alert{(2)} En $\mathbb{R}\times\mathbb{R}=\mathbb{R}^2$,
        tomemos $v_1=v_2=1$, $w_1=0$ y $w_2=1$. Los vectores
        \[
            (v_1,w_1)=(1,0),
            \qquad
            (v_2,w_2)=(1,1)
        \]
        son linealmente independientes, pues
        $\rg\begin{pmatrix}1&1\\0&1\end{pmatrix}=2$.
        Sin embargo, $v_1=v_2$, así que $v_1-v_2=0$ es una relación
        de dependencia lineal entre $v_1,v_2$ en $\mathbb{R}$.

        \framebreak

        \alert{(3)} No. En $\mathbb{R}\times\mathbb{R}=\mathbb{R}^2$,
        tomemos $r=1$ y $v_1=w_1=1$. Entonces $v_1$ genera $\mathbb{R}$
        y $w_1$ genera $\mathbb{R}$, pero el único vector
        \[
            (v_1,w_1)=(1,1)
        \]
        no puede generar $\mathbb{R}^2$, ya que para generar $k^n$
        hacen falta al menos $n$ vectores. Explícitamente, las
        combinaciones lineales de $(1,1)$ son los vectores $(\lambda,\lambda)$,
        y $(1,0)$ no es de esa forma.

        \medskip

        Lo que sí genera $V\times W$ es la familia
        \[
            (v_1,0),\ldots,(v_r,0),\,(0,w_1),\ldots,(0,w_r),
        \]
        como muestra la primera parte de la demostración del teorema
        sobre bases del producto.
    \end{solution}
\end{frame}
